% Part: applied-modal-logics
% Chapter: epistemic-logic
% Section: language-epistemic-logic

\documentclass[../../../include/open-logic-section]{subfiles}

\begin{document}

\olfileid[pt]{aml}{el}{lan}

\olsection{A linguagem da lógica epistêmica}

\begin{defn}
Seja $G$ um conjunto de símbolos de agentes. A linguagem básica da
lógica epistêmica multiagente contém:
\begin{enumerate}
  \tagitem{prvFalse}{A constante proposicional de !!{falsity}~$\lfalse$.}{}
  \tagitem{prvTrue}{A constante proposicional de !!{truth}~$\ltrue$.}{}
  \item Um conjunto de elementos !!{denumerable}s, as !!{propositional variable}s: $\Obj
    p_0$, $\Obj p_1$, $\Obj p_2$, \dots
  \item Os conectivos proposicionais: \startycommalist
  \iftag{prvNot}{\ycomma $\lnot$ (negação)}{}%
  \iftag{prvAnd}{\ycomma $\land$ (conjunção)}{}%
  \iftag{prvOr}{\ycomma $\lor$ (disjunção)}{}%
  \iftag{prvIf}{\ycomma $\lif$ (!!{conditional})}{}%
  \iftag{prvIff}{\ycomma $\liff$ (!!{biconditional})}{}.
  \item O operador de conhecimento $\Knows_a$, em que $a \in G$.
\end{enumerate}
\end{defn}

Se estivermos interessados apenas no conhecimento de um único agente
em nosso sistema, podemos omitir a referência ao conjunto~$G$ e aos
agentes individuais. Nesse caso, temos apenas o operador
básico~$\Knows$.

\begin{defn}
As \emph{!!^{formula}s} da linguagem epistêmica são definidas
  indutivamente da seguinte maneira:
\begin{enumerate}
\tagitem{prvFalse}{$\lfalse$ é uma !!{formula} atômica.}{}

\tagitem{prvTrue}{$\ltrue$ é uma !!{formula} atômica.}{}

\item Toda variável proposicional $\Obj p_i$ é uma !!{formula} (atômica).

\tagitem{prvNot}{Se $!A$ é !!a{formula}, então $\lnot !A$ é
  !!a{formula}.}{}

\tagitem{prvAnd}{Se $!A$ e $!B$ são !!{formula}s, então $(!A \land
  !B)$ é !!a{formula}.}{}

\tagitem{prvOr}{Se $!A$ e $!B$ são !!{formula}s, então $(!A \lor !B)$
  é !!a{formula}.}{}

\tagitem{prvIf}{Se $!A$ e $!B$ são !!{formula}s, então $(!A \lif !B)$
  é !!a{formula}.}{}

\tagitem{prvIff}{Se $!A$ e $!B$ são !!{formula}s, então $(!A \liff !B)$
  é !!a{formula}.}{}

\item Se $!A$ é !!a{formula} e $a \in G$, então $\Knows_a !A$ é
  !!a{formula}.

\tagitem{limitClause}{Nada mais é !!a{formula}.}{}
\end{enumerate}
\end{defn}

Se !!a{formula}~$!A$ não contém $\Knows_a$, dizemos que ela é
\emph{livre de modalidade}.

\begin{defn}
  Enquanto o operador $\Knows$ simboliza o conhecimento individual,
  $\EKnows$, frequentemente lido como ``todos sabem'', simboliza o
  conhecimento do grupo. Se $G' \subseteq G$, definimos $\EKnows_{G'} !A$
 como abreviação de \[\bigwedge_{b \in G'} \Knows_b !A.\]
\end{defn}

Também podemos definir um sentido ainda mais forte de conhecimento: o
\emph{conhecimento comum} entre um grupo de agentes~$G$. Quando uma
informação é conhecimento comum entre um grupo de agentes, isso
significa que, para toda combinação de agentes do grupo, todos sabem
que os demais sabem que os demais sabem \dots ad infinitum. Essa noção
é consideravelmente mais forte do que o conhecimento do grupo, e é
fácil construir modelos relacionais nos quais !!a{formula} é
conhecimento do grupo, mas não conhecimento comum. Usaremos
$\CKnows_G !A$ para simbolizar ``é conhecimento comum entre~$G$
que~$!A$.''

\end{document}
